\section{Rays of Bounded Period}
\label{sec:rays}

After stopping, each moving coordinate can continue indefinitely. We need
continuations that fit a finite description. We let the moving words grow
far beyond their ends (\cref{lem:stopping}), find a short loop in this
extension beyond the end of each moving word, and repeat the loop forever.
Repetition must preserve safety against the stopped targets. A host is
\emph{safe against} a target if no comparison between them carries an
admitted pair.

\subsection{Free windows and free loops}

A loop can change the words to the right of its cuts. To control those
changes, fix a length $\theta>0$ and look at the suffixes of stopped words
of length at most $\theta$, the \emph{short suffixes}. The \emph{right
context} of a cut is the set of short suffixes and the empty word that are
prefixes of the rest at the cut. All members are prefixes of a longest
member, which therefore determines the context. Counting one candidate per
word and length, there are at most $1+k(\theta+1)$ candidates for that
longest member.

A long match against a stopped word ends with one of its suffixes of
length $\theta$. We look for a region where none of these suffixes starts.
Which length works is not known in advance, so we try lengths of
increasing scale, and the depth bound limits how many scales can fail.
For a function $g$, a \emph{free window} at scale $\theta$ consists of
$g(\theta)+1$ consecutive cuts of the host with no such start. The scales
begin at $1$, and the next scale after $\theta$ is $g(\theta)+\theta+1$.
Suppose that scale $\theta$ offers no free window in a region of the host
whose length is the next scale. Then one of the first $g(\theta)+1$ cuts
of the region starts an occurrence of a stopped suffix of length $\theta$.
This occurrence ends strictly before the end of the region, because of the
extra letter in the next scale. Start with a region at the largest scale,
and repeat this step inside each occurrence found, down to scale $1$. If no
scale offers a free window, this gives nested occurrences of stopped
suffixes, one per scale, whose right ends strictly decrease. Trimming them
all to the start of the innermost one gives suffixes of stopped words of
distinct lengths that are prefixes of one rest of the host. A depth bound
of $C$ therefore guarantees a free window within the first $C+1$ scales,
provided the host has room for the next scale.

In a free window, two cuts can repeat both the reader state and the
right context. We call these cuts a \emph{free loop}. The choice
$g(\theta)=N(1+k(\theta+1))$ supplies more cuts than possible pairs. The
interval between the cuts has no start of a stopped suffix of length
$\theta$. The free-window argument and this count give a loop of length at
most
\begin{equation}
 D=(2(k+2)(N+2))^{C+1}.
 \label{eq:period-bound}
\end{equation}

\subsection{Repeating the loop}

Write the extended host as $h=s\ell\tau$, where $\ell$ is the free loop
and $\tau$ is the rest of the host after it. The loop and the word $s$
before it give an infinite word $s\ell^\omega$, called a \emph{ray}. Every
prefix of the ray is a prefix of some $s\ell^n$, and a comparison in a
prefix is also a comparison in the longer word, with the same states. So
it suffices that every $s\ell^n$ is safe against the stopped targets.

The ray has no $\tau$, but $\tau$ serves as a reference for the letters
to the right. Each cut of $s\ell^n$ corresponds to a cut of $h$: a cut in
$s$ to itself, a cut in a copy of $\ell$ to the cut at the same place in
the loop of $h$, and the last cut to the start of the loop. Corresponding cuts have the same state, because
reading $\ell$ returns to its entry state. They also have the same right
context, if the contexts in $s\ell^n$ are computed in $s\ell^n\tau$.
Indeed, the two loop cuts have the same right context, so $\ell\tau$ and
$\tau$ have the same context. Because short suffixes are closed under
taking suffixes, two words with the same right context keep the same
right context when the same word is prepended to both. By induction,
$\ell^n\tau$ has the context of $\tau$.

Now let a stopped target rest $y$ fit at a cut of $s\ell^n$. If
$\len y\le\theta$, then $y$ is in the right context, so it fits at the
corresponding cut of $h$. If $\len y>\theta$, its last $\theta$ letters
form a short suffix that starts at a cut of $s\ell^n$. That cut cannot lie
in a copy of $\ell$, since the corresponding cut of $h$ would start a
suffix of length $\theta$ inside the free window. So it lies in $s$, where
the earlier letters of the match are unchanged, and $y$ fits at the same
cut of $h$. In both cases the comparison occurs in $h$ with the same pair
of states. So the ray inherits safety from $h$, and the repeated loop keeps
the run alive.

\begin{lemma}[Periodic family]\label{lem:periodic}
From the stopped extension of \cref{lem:stopping} with depth at most $C$,
one can choose a ray of period at most $D$ for each moving coordinate.
The ray starts with the word at that coordinate. For every $e\ge 0$,
cutting all rays at the lengths of their starting words plus $e$ gives a
hole together with the stopped words.
\end{lemma}

The proof takes a free loop after the end of each moving word, using a
sufficiently long stage of joint extension. Rejection persists and every
prefix is safe against the stopped targets. By \cref{lem:joint}, these
properties allow the rays to advance together. The details of the free
windows, the period estimate, and the repetition argument are in
\cref{app:rays}. In \cref{fig:stopping}, the moving word $0$ gives the ray
$0^\omega$ with period $1$.

\subsection{Offsets between rays}
\label{sec:offsets}

The rays grow in step, and their common truncations must avoid admitted
comparisons between rays. Cut every ray $i$ at $o_i+e$ for a common
increment $e$. Suppose the infinite rests at a host cut $a$ of ray $i$ and
a target cut $b$ of ray $j$ are equal. Once $e$ is large enough to reach
both cuts, their finite rests are prefixes of one infinite word, of lengths
$o_i+e-a$ and $o_j+e-b$. They form a comparison exactly when
$o_j+e-b\le o_i+e-a$, which does not depend on $e$. If the pair of states
is admitted, the comparison must be excluded:
\begin{equation}
   o_i+b-a+1\le o_j .
   \label{eq:offset}
\end{equation}
We call the stopped words and moving rays a \emph{periodic tuple} if the
stopped words are eligible, the rays are live and eventually rejected,
and the stopped words and every ray prefix are safe against each stopped
target. In addition it has integer \emph{offsets} $o_i$ such that every
pair of moving cuts with equal infinite rests and an admitted pair of
states satisfies \eqref{eq:offset}. The condition ranges over all ordered
pairs of moving coordinates and all their cuts. By the computation above,
the starting lengths in \cref{lem:periodic} are offsets. The offsets link
the next two steps: the trie descent keeps them, and the truncation turns
them into lengths.

The period bound leaves the lengths of the finite parts uncontrolled.
The next step shortens these parts without losing the periodic tuple.
